\subsubsection{Min-Cost Max-Flow}
\begin{lstlisting}[style=mystyle, language=C++]
// TC: O(F * N * M) with SPFA; swap for Dijkstra + potentials to get O(F * (N+M) log N)
// usa: maximo flujo con costo minimo (asume costos no negativos)
#include <bits/stdc++.h>
using namespace std;

const long long INF = 1e18;

struct MCMF{
	// to = destino, rev = indice de la arista inversa en g[to]
	// cap = capacidad restante, cost = costo unitario
	struct Edge{ int to, rev; long long cap, cost; };
	int n;
	// dist = distancia reducida desde s
	vector<vector<Edge>> g;
	// pot = potentials para Dijkstra (h(v) en reweighting)
	vector<long long> dist, pot;
	// parent_v = nodo anterior en el camino aumentante, parent_e = indice de la arista usada
	vector<int> parent_v, parent_e;
	MCMF(int n): n(n){
		g.assign(n, {});
		dist.assign(n, INF);
		pot.assign(n, 0);
		parent_v.assign(n, -1);
		parent_e.assign(n, -1);
	}
	void addEdge(int u, int v, long long cap, long long cost){
		Edge a{v, (int)g[v].size(), cap, cost};   // arista real
		Edge b{u, (int)g[u].size(), 0, -cost};    // arista residual (cap 0, costo negativo)
		g[u].push_back(a);
		g[v].push_back(b);
	}
	// returns (maxflow, mincost)
	pair<long long,long long> minCostMaxFlow(int s, int t){
		long long flow = 0, cost = 0;
		while(true){
			// SPFA con potentials: busca el camino de costo minimo de s a t en el grafo residual
			fill(dist.begin(), dist.end(), INF);
			dist[s] = 0;
			vector<bool> inq(n, false);       // inq = nodo actualmente en la cola de SPFA
			queue<int> q;
			q.push(s); inq[s] = true;
			while(!q.empty()){
				int v = q.front(); q.pop();
				inq[v] = false;
				for(int i=0;i<(int)g[v].size();i++){
					auto& e = g[v][i];
					// arista relajable con costos reducidos (pot[v] - pot[e.to] >= 0)
					if(e.cap > 0 && dist[v] + e.cost + pot[v] - pot[e.to] < dist[e.to]){
						dist[e.to] = dist[v] + e.cost + pot[v] - pot[e.to];
						parent_v[e.to] = v;
						parent_e[e.to] = i;
						if(!inq[e.to]){ inq[e.to] = true; q.push(e.to); }
					}
				}
			}
			if(dist[t] == INF) break;   // no hay mas camino de s a t
			// actualiza potentials para mantener costos reducidos no negativos
			for(int v=0;v<n;v++) if(dist[v] < INF) pot[v] += dist[v];
			// encuentra el cuello de botella (minima capacidad en el camino)
			long long add = INF;
			for(int v=t; v != s; v = parent_v[v]){
				add = min(add, g[parent_v[v]][parent_e[v]].cap);
			}
			// empuja flujo por el camino y suma el costo
			for(int v=t; v != s; v = parent_v[v]){
				auto& e = g[parent_v[v]][parent_e[v]];
				e.cap -= add;
				g[v][e.rev].cap += add;
				cost += add * e.cost;
			}
			flow += add;
		}
		return {flow, cost};
	}
};

int main(){
	MCMF mcmf(4);
	mcmf.addEdge(0,1,10,1);
	mcmf.addEdge(0,2,5,2);
	mcmf.addEdge(1,3,10,1);
	mcmf.addEdge(2,3,10,1);
	auto [flow, cost] = mcmf.minCostMaxFlow(0, 3);
	cout << "flow=" << flow << " cost=" << cost << "\n";
	return 0;
}
\end{lstlisting}
